\documentclass[aspectratio=169,11pt]{beamer}

% ============================================================
% CLASE 6 - ESTADÍSTICA PARA DATA SCIENCE
% Diplomado de Python y Análisis de Datos
% Universidad Marista
%
% Compatible con:
% Overleaf + pdfLaTeX
% ============================================================

% ============================================================
% CODIFICACIÓN
% ============================================================

\usepackage[utf8]{inputenc}
\usepackage[T1]{fontenc}
\usepackage[spanish,es-nodecimaldot]{babel}
\usepackage{lmodern}

% ============================================================
% PAQUETES
% ============================================================

\usepackage{amsmath}
\usepackage{amssymb}
\usepackage{booktabs}
\usepackage{array}
\usepackage{multicol}
\usepackage{graphicx}
\usepackage{xcolor}
\usepackage{tikz}
\usepackage{listings}

\usetikzlibrary{
    positioning,
    arrows.meta,
    shapes.geometric,
    calc
}

% ============================================================
% TEMA BEAMER
% ============================================================

\usetheme{Madrid}
\usecolortheme{default}

\setbeamertemplate{navigation symbols}{}

% ============================================================
% COLORES
% ============================================================

\definecolor{azulmarista}{RGB}{18,67,115}
\definecolor{azulclaro}{RGB}{224,238,250}

\definecolor{verde}{RGB}{38,120,85}
\definecolor{verdeclaro}{RGB}{225,245,235}

\definecolor{naranja}{RGB}{225,125,35}
\definecolor{naranjaclaro}{RGB}{253,237,218}

\definecolor{morado}{RGB}{104,72,150}
\definecolor{moradoclaro}{RGB}{238,230,248}

\definecolor{gris}{RGB}{80,80,80}
\definecolor{grisclaro}{RGB}{245,245,245}

\definecolor{rojo}{RGB}{175,50,50}
\definecolor{rojoclaro}{RGB}{250,230,230}

% ============================================================
% COLORES BEAMER
% ============================================================

\setbeamercolor{structure}{
    fg=azulmarista
}

\setbeamercolor{frametitle}{
    fg=white,
    bg=azulmarista
}

\setbeamercolor{title}{
    fg=azulmarista
}

\setbeamercolor{subtitle}{
    fg=gris
}

\setbeamercolor{block title}{
    fg=white,
    bg=azulmarista
}

\setbeamercolor{block body}{
    bg=azulclaro
}

\setbeamercolor{block title example}{
    fg=white,
    bg=verde
}

\setbeamercolor{block body example}{
    bg=verdeclaro
}

\setbeamercolor{block title alerted}{
    fg=white,
    bg=rojo
}

\setbeamercolor{block body alerted}{
    bg=rojoclaro
}

% ============================================================
% PIE DE PÁGINA
% ============================================================

\setbeamertemplate{footline}{
    \leavevmode
    \hbox{
        \begin{beamercolorbox}[
            wd=.74\paperwidth,
            ht=2.7ex,
            dp=1.1ex,
            leftskip=1em
        ]{author in head/foot}
            Diplomado de Python y Análisis de Datos
        \end{beamercolorbox}
        \begin{beamercolorbox}[
            wd=.26\paperwidth,
            ht=2.7ex,
            dp=1.1ex,
            center
        ]{date in head/foot}
            \insertframenumber/\inserttotalframenumber
        \end{beamercolorbox}
    }
}

% ============================================================
% CÓDIGO PYTHON
% ============================================================

\lstdefinestyle{pythonstyle}{
    language=Python,
    backgroundcolor=\color{grisclaro},
    basicstyle=\ttfamily\scriptsize,
    keywordstyle=\color{azulmarista}\bfseries,
    commentstyle=\color{verde},
    stringstyle=\color{morado},
    numberstyle=\tiny\color{gris},
    numbers=left,
    numbersep=5pt,
    breaklines=true,
    breakatwhitespace=false,
    keepspaces=true,
    showspaces=false,
    showstringspaces=false,
    showtabs=false,
    tabsize=4,
    frame=single,
    rulecolor=\color{azulmarista}
}

\lstset{
    style=pythonstyle
}

% ============================================================
% INFORMACIÓN
% ============================================================

\title{Estadística para Data Science}

\subtitle{
Clase 6\\
De describir datos a cuantificar incertidumbre
}

\author{
Diplomado de Python y Análisis de Datos
}

\institute{
Universidad Marista
}

\date{2026}

% ============================================================
% DOCUMENTO
% ============================================================

\begin{document}

% ============================================================
% PORTADA
% ============================================================

\begin{frame}
    \titlepage
\end{frame}

% ============================================================
% OBJETIVO GENERAL
% ============================================================

\begin{frame}{Objetivo de la sesión}

Al finalizar la clase, el estudiante será capaz de:

\begin{itemize}

    \item distinguir entre población y muestra;

    \item diferenciar parámetros y estadísticos;

    \item interpretar medidas de tendencia central y dispersión;

    \item analizar distribuciones;

    \item comprender el Teorema Central del Límite;

    \item interpretar el error estándar;

    \item construir intervalos de confianza;

    \item formular hipótesis estadísticas;

    \item interpretar correctamente un valor $p$;

    \item comparar dos grupos mediante una prueba t;

    \item distinguir significancia estadística de relevancia práctica.

\end{itemize}

\end{frame}

% ============================================================
% CRONOGRAMA
% ============================================================

\begin{frame}{Ruta de las 3 horas}

\scriptsize

\begin{table}
\centering

\begin{tabular}{p{2.2cm}p{8.7cm}}
\toprule

\textbf{Tiempo} &
\textbf{Tema} \\

\midrule

0:00--0:10 &
Problema inicial: ¿una diferencia observada es real? \\

0:10--0:25 &
Población, muestra, parámetros y estadísticos \\

0:25--0:45 &
Media, mediana, varianza, desviación y percentiles \\

0:45--1:00 &
Distribuciones y distribución normal \\

1:00--1:20 &
Muestreo y Teorema Central del Límite \\

1:20--1:30 &
Error estándar e introducción a inferencia \\

1:30--1:40 &
Descanso \\

1:40--1:55 &
Intervalos de confianza \\

1:55--2:15 &
Hipótesis nula, alternativa y valor $p$ \\

2:15--2:30 &
Prueba t y comparación de medias \\

2:30--2:40 &
Correlación y significancia práctica \\

2:40--2:55 &
Statistical Detective Challenge \\

2:55--3:00 &
Conclusiones y conexión con Series Temporales \\

\bottomrule

\end{tabular}

\end{table}

\end{frame}

% ============================================================
\section{Introducción}
% ============================================================

\begin{frame}{El problema inicial}

Supongamos que una empresa prueba dos versiones de una página web.

\vspace{0.3cm}

\begin{center}

\begin{tabular}{lcc}
\toprule

&
\textbf{Versión A}
&
\textbf{Versión B} \\

\midrule

Conversión &
12.1\% &
13.0\% \\

\bottomrule
\end{tabular}

\end{center}

\vspace{0.5cm}

\begin{alertblock}{Pregunta}
¿Podemos afirmar inmediatamente que la versión B es mejor?
\end{alertblock}

\vspace{0.3cm}

No necesariamente.

\end{frame}

% ------------------------------------------------------------

\begin{frame}{¿Qué nos falta saber?}

Una diferencia observada puede depender de:

\begin{itemize}

    \item tamaño de muestra;

    \item variabilidad de los datos;

    \item forma de la distribución;

    \item diseño del experimento;

    \item sesgo de selección;

    \item incertidumbre muestral;

    \item magnitud real del efecto.

\end{itemize}

\vspace{0.4cm}

\begin{block}{Idea central}
La estadística nos ayuda a distinguir entre señal y ruido.
\end{block}

\end{frame}

% ------------------------------------------------------------

\begin{frame}{De EDA a inferencia}

\begin{center}

\begin{tikzpicture}[
node distance=0.7cm,
every node/.style={
draw,
rounded corners,
fill=azulclaro,
minimum width=1.8cm,
minimum height=0.7cm,
font=\small
}
]

\node (datos) {Datos};

\node[right=of datos] (eda) {EDA};

\node[right=of eda] (patron) {Patrón};

\node[right=of patron] (estadistica) {Estadística};

\node[below=1cm of estadistica] (evidencia) {Evidencia};

\node[left=of evidencia] (decision) {Decisión};

\draw[-{Latex}] (datos) -- (eda);

\draw[-{Latex}] (eda) -- (patron);

\draw[-{Latex}] (patron) -- (estadistica);

\draw[-{Latex}] (estadistica) -- (evidencia);

\draw[-{Latex}] (evidencia) -- (decision);

\end{tikzpicture}

\end{center}

\vspace{0.4cm}

\begin{block}{Pregunta}
¿El patrón que observamos es suficientemente sólido para tomar una decisión?
\end{block}

\end{frame}

% ============================================================
\section{Descriptiva e inferencial}
% ============================================================

\begin{frame}{Estadística descriptiva}

La estadística descriptiva resume los datos observados.

\vspace{0.3cm}

Herramientas:

\begin{itemize}

    \item media;

    \item mediana;

    \item moda;

    \item varianza;

    \item desviación estándar;

    \item percentiles;

    \item tablas;

    \item visualizaciones.

\end{itemize}

\vspace{0.4cm}

\begin{exampleblock}{Ejemplo}
¿Cuál fue el gasto promedio de los clientes observados?
\end{exampleblock}

\end{frame}

% ------------------------------------------------------------

\begin{frame}{Estadística inferencial}

La estadística inferencial busca utilizar una muestra para obtener
conclusiones sobre una población.

\vspace{0.4cm}

Herramientas:

\begin{itemize}

    \item estimación;

    \item distribuciones muestrales;

    \item error estándar;

    \item intervalos de confianza;

    \item pruebas de hipótesis;

    \item significancia estadística.

\end{itemize}

\vspace{0.4cm}

\begin{exampleblock}{Ejemplo}
¿Podemos inferir el gasto promedio de todos los clientes utilizando una muestra?
\end{exampleblock}

\end{frame}

% ------------------------------------------------------------

\begin{frame}{Descriptiva vs inferencial}

\begin{columns}

\column{0.45\textwidth}

\begin{block}{Descriptiva}

Pregunta:

\begin{center}
¿Qué ocurrió en los datos observados?
\end{center}

\end{block}

\column{0.45\textwidth}

\begin{exampleblock}{Inferencial}

Pregunta:

\begin{center}
¿Qué podemos concluir más allá de los datos observados?
\end{center}

\end{exampleblock}

\end{columns}

\end{frame}

% ============================================================
\section{Población y muestra}
% ============================================================

\begin{frame}{¿Qué es una población?}

Una \textbf{población} es el conjunto completo de elementos que queremos estudiar.

Ejemplos:

\begin{itemize}

    \item todos los clientes de una empresa;

    \item todos los estudiantes de una universidad;

    \item todas las transacciones de un año;

    \item todas las piezas producidas por una fábrica;

    \item todos los usuarios de una aplicación.

\end{itemize}

\end{frame}

% ------------------------------------------------------------

\begin{frame}{¿Qué es una muestra?}

Una \textbf{muestra} es un subconjunto de la población.

\vspace{0.3cm}

\begin{center}

\begin{tikzpicture}

\draw[
    fill=azulclaro,
    draw=azulmarista,
    thick
]
(0,0) circle (2);

\node at (0,1.35) {
    \textbf{Población}
};

\draw[
    fill=verdeclaro,
    draw=verde,
    thick
]
(0,-0.4) circle (0.9);

\node at (0,-0.4) {
    \textbf{Muestra}
};

\end{tikzpicture}

\end{center}

\end{frame}

% ------------------------------------------------------------

\begin{frame}{¿Por qué trabajar con muestras?}

Porque estudiar toda la población puede ser:

\begin{itemize}

    \item muy costoso;

    \item muy lento;

    \item computacionalmente complejo;

    \item imposible;

    \item innecesario cuando una muestra adecuada puede proporcionar una buena estimación.

\end{itemize}

\end{frame}

% ------------------------------------------------------------

\begin{frame}{Parámetro vs estadístico}

\begin{columns}

\column{0.45\textwidth}

\begin{block}{Parámetro}

Describe a la población.

Ejemplo:

\[
\mu
\]

Media poblacional.

\end{block}

\column{0.45\textwidth}

\begin{exampleblock}{Estadístico}

Se calcula con una muestra.

Ejemplo:

\[
\bar{x}
\]

Media muestral.

\end{exampleblock}

\end{columns}

\vspace{0.5cm}

\begin{center}

Muestra

$\downarrow$

Estadístico

$\downarrow$

Estimación del parámetro

\end{center}

\end{frame}

% ============================================================
\section{Sesgo de muestreo}
% ============================================================

\begin{frame}{Una muestra puede estar sesgada}

Supongamos:

\begin{quote}

Queremos medir la satisfacción de todos los estudiantes,
pero solamente encuestamos a quienes asistieron hoy.

\end{quote}

\vspace{0.5cm}

\begin{alertblock}{Pregunta}
¿La muestra representa necesariamente a toda la población?
\end{alertblock}

No.

\end{frame}

% ------------------------------------------------------------

\begin{frame}{Ejemplo de sesgo}

Queremos conocer el uso de redes sociales de una población.

Pero publicamos la encuesta solamente en Instagram.

\vspace{0.4cm}

La muestra podría sobre-representar:

\begin{itemize}

    \item usuarios activos de Instagram;

    \item personas con determinadas edades;

    \item ciertos hábitos digitales.

\end{itemize}

\vspace{0.3cm}

\begin{block}{Idea}
El proceso de selección de la muestra importa.
\end{block}

\end{frame}

% ------------------------------------------------------------

\begin{frame}{Más datos no corrigen todo}

Podemos tener:

\begin{center}

{\Huge
1\,000\,000
}

\vspace{0.3cm}

observaciones.

\end{center}

Pero si la muestra está sistemáticamente sesgada:

\begin{center}

\Large
más datos $\neq$ mejor inferencia
\end{center}

\vspace{0.4cm}

\begin{alertblock}{Lección}
Una muestra enorme puede seguir representando mal a la población.
\end{alertblock}

\end{frame}

% ============================================================
\section{Tendencia central}
% ============================================================

\begin{frame}[fragile]{Media}

La media aritmética es:

\[
\bar{x}
=
\frac{1}{n}
\sum_{i=1}^{n}x_i
\]

En Pandas:

\begin{lstlisting}
media = df["ingreso"].mean()

print(media)
\end{lstlisting}

\begin{block}{Interpretación}
Representa el centro aritmético de las observaciones.
\end{block}

\end{frame}

% ------------------------------------------------------------

\begin{frame}[fragile]{Mediana}

La mediana es el valor central después de ordenar las observaciones.

\begin{lstlisting}
mediana = df["ingreso"].median()

print(mediana)
\end{lstlisting}

Ejemplo:

\begin{center}

10,\ 11,\ 12,\ 13,\ 1000

\end{center}

\end{frame}

% ------------------------------------------------------------

\begin{frame}{Media vs mediana}

Para:

\[
10,\ 11,\ 12,\ 13,\ 1000
\]

tenemos:

\begin{columns}

\column{0.45\textwidth}

\begin{block}{Media}

\[
209.2
\]

\end{block}

\column{0.45\textwidth}

\begin{exampleblock}{Mediana}

\[
12
\]

\end{exampleblock}

\end{columns}

\vspace{0.5cm}

\begin{alertblock}{Pregunta}
¿Cuál representa mejor el valor típico de este grupo?
\end{alertblock}

\end{frame}

% ------------------------------------------------------------

\begin{frame}[fragile]{Moda}

La moda representa el valor más frecuente.

\begin{lstlisting}
moda = df["categoria"].mode()

print(moda)
\end{lstlisting}

Es muy útil para variables categóricas.

Ejemplos:

\begin{itemize}

    \item categoría más frecuente;

    \item ciudad más común;

    \item producto más observado;

    \item tipo de membresía dominante.

\end{itemize}

\end{frame}

% ============================================================
\section{Dispersión}
% ============================================================

\begin{frame}{El promedio no cuenta toda la historia}

Consideremos:

\vspace{0.2cm}

\begin{columns}

\column{0.45\textwidth}

\begin{block}{Grupo A}

\[
48,\ 49,\ 50,\ 51,\ 52
\]

\end{block}

\column{0.45\textwidth}

\begin{exampleblock}{Grupo B}

\[
10,\ 30,\ 50,\ 70,\ 90
\]

\end{exampleblock}

\end{columns}

\vspace{0.5cm}

Ambos tienen media:

\[
50
\]

pero presentan comportamientos muy diferentes.

\end{frame}

% ------------------------------------------------------------

\begin{frame}[fragile]{Rango}

El rango mide la distancia entre el máximo y el mínimo:

\[
R
=
x_{\max}
-
x_{\min}
\]

En Pandas:

\begin{lstlisting}
rango = (
    df["ventas"].max()
    - df["ventas"].min()
)

print(rango)
\end{lstlisting}

\begin{alertblock}{Limitación}
El rango depende solamente del mínimo y máximo.
\end{alertblock}

\end{frame}

% ------------------------------------------------------------

\begin{frame}[fragile]{Varianza}

La varianza muestral puede expresarse como:

\[
s^2
=
\frac{
\sum_{i=1}^{n}
(x_i-\bar{x})^2
}{
n-1
}
\]

En Pandas:

\begin{lstlisting}
varianza = df["ventas"].var()

print(varianza)
\end{lstlisting}

\begin{block}{Interpretación}
Mide cuánto se alejan, en promedio cuadrático, los datos respecto de la media.
\end{block}

\end{frame}

% ------------------------------------------------------------

\begin{frame}[fragile]{Desviación estándar}

La desviación estándar es:

\[
s
=
\sqrt{s^2}
\]

En Pandas:

\begin{lstlisting}
desviacion = df["ventas"].std()

print(desviacion)
\end{lstlisting}

\vspace{0.3cm}

\begin{block}{Idea}
Una desviación estándar mayor indica mayor dispersión.
\end{block}

\end{frame}

% ------------------------------------------------------------

\begin{frame}{Misma media, diferente dispersión}

\begin{center}

\begin{tabular}{lcc}
\toprule

&
\textbf{Grupo A}
&
\textbf{Grupo B} \\

\midrule

Media &
50 &
50 \\

Dispersión &
Baja &
Alta \\

\bottomrule

\end{tabular}

\end{center}

\vspace{0.5cm}

\begin{alertblock}{Lección}
Una media debería interpretarse junto con una medida de dispersión.
\end{alertblock}

\end{frame}

% ------------------------------------------------------------

\begin{frame}{Coeficiente de variación}

Una medida relativa de dispersión es:

\[
CV
=
\frac{s}{\bar{x}}
\]

Ejemplo:

\begin{center}

\begin{tabular}{lcc}
\toprule

&
Producto A &
Producto B \\

\midrule

Media &
100 &
1000 \\

Desv. estándar &
10 &
50 \\

\bottomrule

\end{tabular}

\end{center}

\vspace{0.4cm}

El coeficiente permite comparar variabilidad relativa.

\end{frame}

% ============================================================
\section{Percentiles}
% ============================================================

\begin{frame}[fragile]{Percentiles}

Los percentiles indican la posición relativa de un valor.

En Pandas:

\begin{lstlisting}
q1 = df["ventas"].quantile(0.25)

mediana = df["ventas"].quantile(0.50)

q3 = df["ventas"].quantile(0.75)

p95 = df["ventas"].quantile(0.95)

print(q1, mediana, q3, p95)
\end{lstlisting}

\end{frame}

% ------------------------------------------------------------

\begin{frame}{Interpretación de percentiles}

Supongamos que un cliente está en el:

\[
P_{95}
\]

de gasto.

Podemos interpretar que su gasto es superior aproximadamente al
95\% de las observaciones del dataset.

\vspace{0.4cm}

Aplicaciones:

\begin{itemize}

    \item segmentación de clientes;

    \item detección de valores extremos;

    \item análisis de rendimiento;

    \item riesgo;

    \item calidad.

\end{itemize}

\end{frame}

% ============================================================
\section{Distribuciones}
% ============================================================

\begin{frame}{¿Qué es una distribución?}

Una distribución describe:

\begin{itemize}

    \item dónde se concentran los datos;

    \item qué tan dispersos están;

    \item qué forma tienen;

    \item si existe asimetría;

    \item si existen colas;

    \item si aparecen valores extremos;

    \item si existen varios grupos.

\end{itemize}

\end{frame}

% ------------------------------------------------------------

\begin{frame}[fragile]{Visualizar una distribución}

\begin{lstlisting}
import seaborn as sns
import matplotlib.pyplot as plt

sns.histplot(
    data=df,
    x="ventas",
    kde=True
)

plt.title(
    "Distribucion de ventas"
)

plt.show()
\end{lstlisting}

\end{frame}

% ------------------------------------------------------------

\begin{frame}{¿Qué buscamos en una distribución?}

\begin{itemize}

    \item simetría;

    \item sesgo positivo;

    \item sesgo negativo;

    \item concentración;

    \item dispersión;

    \item multimodalidad;

    \item valores extremos.

\end{itemize}

\vspace{0.4cm}

\begin{block}{Pregunta}
¿Qué información observamos aquí que una media no puede mostrarnos?
\end{block}

\end{frame}

% ============================================================
\section{Distribución normal}
% ============================================================

\begin{frame}{Distribución normal}

La distribución normal es uno de los modelos probabilísticos más importantes.

Características:

\begin{itemize}

    \item simétrica;

    \item forma de campana;

    \item máxima concentración cerca del centro;

    \item menor frecuencia hacia las colas;

    \item media, mediana y moda coinciden en el modelo ideal.

\end{itemize}

\begin{alertblock}{Importante}
No todos los datos reales siguen una distribución normal.
\end{alertblock}

\end{frame}

% ------------------------------------------------------------

\begin{frame}{Regla 68--95--99.7}

Para una variable aproximadamente normal:

\[
\mu \pm 1\sigma
\approx
68\%
\]

\[
\mu \pm 2\sigma
\approx
95\%
\]

\[
\mu \pm 3\sigma
\approx
99.7\%
\]

\end{frame}

% ------------------------------------------------------------

\begin{frame}{Ejemplo de la regla empírica}

Supongamos:

\[
\mu = 100
\]

\[
\sigma = 10
\]

Entonces aproximadamente:

\[
68\%
\rightarrow
[90,110]
\]

\[
95\%
\rightarrow
[80,120]
\]

\[
99.7\%
\rightarrow
[70,130]
\]

\end{frame}

% ============================================================
\section{Z-score}
% ============================================================

\begin{frame}{Z-score}

El z-score indica cuántas desviaciones estándar separan un valor de la media:

\[
z
=
\frac{x-\mu}{\sigma}
\]

Ejemplo:

\[
\mu=100
\]

\[
\sigma=10
\]

\[
x=130
\]

Entonces:

\[
z=3
\]

\end{frame}

% ------------------------------------------------------------

\begin{frame}[fragile]{Z-score en Python}

\begin{lstlisting}
from scipy.stats import zscore

df["z_ventas"] = zscore(
    df["ventas"]
)

print(
    df[
        ["ventas", "z_ventas"]
    ].head()
)
\end{lstlisting}

Aplicaciones:

\begin{itemize}

    \item estandarización;

    \item detección exploratoria de extremos;

    \item comparación entre escalas.

\end{itemize}

\end{frame}

% ============================================================
\section{Distribuciones sesgadas}
% ============================================================

\begin{frame}{Ejemplo de sesgo positivo}

Consideremos ingresos:

\[
15000,\ 20000,\ 23000,\ 27000,\ 30000,\ 300000
\]

La mayoría está relativamente concentrada.

Pero existe un valor muy grande.

\vspace{0.3cm}

Esto genera una cola hacia la derecha.

\vspace{0.3cm}

Frecuentemente:

\[
\text{media}
>
\text{mediana}
\]

\end{frame}

% ------------------------------------------------------------

\begin{frame}{¿Media o mediana?}

\begin{block}{Media}
Útil cuando queremos representar el centro aritmético y los extremos no dominan excesivamente.
\end{block}

\vspace{0.3cm}

\begin{exampleblock}{Mediana}
Especialmente útil cuando existen valores extremos o distribuciones fuertemente sesgadas.
\end{exampleblock}

\vspace{0.3cm}

\begin{alertblock}{Lección}
No existe una medida universalmente mejor.
Depende del problema y de la distribución.
\end{alertblock}

\end{frame}

% ============================================================
\section{Muestreo}
% ============================================================

\begin{frame}{Variabilidad muestral}

Supongamos que tomamos diferentes muestras de la misma población.

\begin{center}

\begin{tabular}{lc}
\toprule

Muestra &
Media \\

\midrule

A &
850 \\

B &
872 \\

C &
844 \\

\bottomrule

\end{tabular}

\end{center}

\vspace{0.4cm}

\begin{alertblock}{Pregunta}
¿Por qué cambia la media?
\end{alertblock}

Porque cada muestra contiene observaciones diferentes.

\end{frame}

% ------------------------------------------------------------

\begin{frame}[fragile]{Tomar una muestra con Pandas}

\begin{lstlisting}
muestra = df.sample(
    n=100,
    random_state=42
)

media = muestra["ventas"].mean()

print(media)
\end{lstlisting}

Si cambiamos la muestra, la media puede cambiar.

\end{frame}

% ------------------------------------------------------------

\begin{frame}[fragile]{Cambiar la muestra}

\begin{lstlisting}
muestra_1 = df.sample(
    n=100,
    random_state=42
)

muestra_2 = df.sample(
    n=100,
    random_state=7
)

print(
    muestra_1["ventas"].mean()
)

print(
    muestra_2["ventas"].mean()
)
\end{lstlisting}

\end{frame}

% ------------------------------------------------------------

\begin{frame}{Tamaño de muestra}

En general, cuando aumenta el tamaño de una muestra aleatoria adecuada:

\begin{itemize}

    \item las estimaciones tienden a ser más estables;

    \item disminuye parte de la incertidumbre muestral;

    \item mejora la precisión estadística.

\end{itemize}

\vspace{0.4cm}

\begin{alertblock}{Pero...}
Aumentar $n$ no corrige automáticamente sesgos sistemáticos.
\end{alertblock}

\end{frame}

% ------------------------------------------------------------

\begin{frame}[fragile]{Comparar tamaños de muestra}

\begin{lstlisting}
m10 = (
    df.sample(10)["ventas"]
    .mean()
)

m100 = (
    df.sample(100)["ventas"]
    .mean()
)

m1000 = (
    df.sample(1000)["ventas"]
    .mean()
)

print(m10)
print(m100)
print(m1000)
\end{lstlisting}

\end{frame}

% ============================================================
\section{Ley de los Grandes Números}
% ============================================================

\begin{frame}{Ley de los Grandes Números}

De manera conceptual:

\begin{quote}

A medida que aumenta el tamaño de una muestra aleatoria adecuada,
la media muestral tiende a acercarse a la media poblacional.

\end{quote}

Ejemplo conceptual:

\begin{center}

\begin{tabular}{lr}
\toprule

Tamaño &
Media \\

\midrule

10 &
930 \\

100 &
861 \\

1000 &
852 \\

10000 &
849 \\

\bottomrule

\end{tabular}

\end{center}

\end{frame}

% ============================================================
\section{Teorema Central del Límite}
% ============================================================

\begin{frame}{Teorema Central del Límite}

Supongamos que tomamos muchas muestras aleatorias de una población.

Para cada muestra calculamos:

\[
\bar{x}
\]

Si repetimos el procedimiento muchas veces, obtenemos una
\textbf{distribución de medias muestrales}.

\vspace{0.3cm}

\begin{block}{Teorema Central del Límite}
Bajo condiciones apropiadas, cuando el tamaño de muestra aumenta,
la distribución de la media muestral tiende aproximadamente a una distribución normal.
\end{block}

\end{frame}

% ------------------------------------------------------------

\begin{frame}{Una idea sorprendente}

La población original:

\begin{center}

\Large
\textbf{no necesita ser normal}

\end{center}

\vspace{0.5cm}

Sin embargo, bajo condiciones apropiadas y con muestras suficientemente grandes:

\[
\bar{X}
\]

puede presentar una distribución aproximadamente normal.

\vspace{0.3cm}

\begin{alertblock}{Importancia}
Esta propiedad es fundamental para muchos procedimientos de inferencia.
\end{alertblock}

\end{frame}

% ------------------------------------------------------------

\begin{frame}{De población a distribución muestral}

\begin{center}

\begin{tikzpicture}[
node distance=0.55cm,
every node/.style={
    draw,
    rounded corners,
    minimum width=3.5cm,
    minimum height=0.65cm,
    align=center,
    font=\small
}
]

\node[
    fill=naranjaclaro
] (pob) {
    Población\\
    posiblemente sesgada
};

\node[
    below=of pob,
    fill=azulclaro
] (muestras) {
    Muchas muestras\\
    de tamaño $n$
};

\node[
    below=of muestras,
    fill=verdeclaro
] (medias) {
    Calculamos\\
    $\bar{x}_1,\bar{x}_2,\ldots,\bar{x}_k$
};

\node[
    below=of medias,
    fill=moradoclaro
] (dist) {
    Distribución de\\
    medias muestrales
};

\draw[-{Latex}] (pob) -- (muestras);

\draw[-{Latex}] (muestras) -- (medias);

\draw[-{Latex}] (medias) -- (dist);

\end{tikzpicture}

\end{center}

\end{frame}

% ------------------------------------------------------------

\begin{frame}[fragile]{Simulando el Teorema Central del Límite}

\begin{lstlisting}
import numpy as np
import matplotlib.pyplot as plt

poblacion = np.random.exponential(
    scale=10,
    size=100000
)

medias = []

for _ in range(5000):

    muestra = np.random.choice(
        poblacion,
        size=50
    )

    medias.append(
        muestra.mean()
    )

plt.hist(
    medias,
    bins=40
)

plt.title(
    "Distribucion de medias muestrales"
)

plt.xlabel("Media")
plt.ylabel("Frecuencia")

plt.show()
\end{lstlisting}

\end{frame}

% ============================================================
\section{Error estándar}
% ============================================================

\begin{frame}{Error estándar}

La desviación estándar describe la dispersión de las observaciones.

El \textbf{error estándar} describe la variabilidad esperada de una estimación.

Para la media:

\[
SE
=
\frac{s}{\sqrt{n}}
\]

donde:

\begin{itemize}

    \item $s$ = desviación estándar muestral;

    \item $n$ = tamaño de muestra.

\end{itemize}

\end{frame}

% ------------------------------------------------------------

\begin{frame}{¿Qué ocurre cuando aumenta $n$?}

Tenemos:

\[
SE
=
\frac{s}{\sqrt{n}}
\]

Si aumenta:

\[
n
\]

entonces:

\[
SE
\downarrow
\]

\vspace{0.4cm}

\begin{block}{Interpretación}
En general, una muestra mayor permite estimar la media con mayor precisión,
si el diseño de muestreo es adecuado.
\end{block}

\end{frame}

% ============================================================
% DESCANSO
% ============================================================

\begin{frame}

\begin{center}

{\Huge
\textcolor{azulmarista}{
\textbf{DESCANSO}
}
}

\vspace{0.5cm}

{\Large
10 minutos
}

\vspace{1cm}

Al regresar:

\vspace{0.3cm}

{\Large
¿Cómo cuantificamos la incertidumbre de una estimación?
}

\end{center}

\end{frame}

% ============================================================
\section{Intervalos de confianza}
% ============================================================

\begin{frame}{Una estimación puntual}

Supongamos que nuestra muestra produce:

\[
\bar{x}=850
\]

Podemos decir:

\begin{center}

\Large
``La media estimada es 850''

\end{center}

\vspace{0.4cm}

Pero otra muestra podría producir:

\[
842
\]

o:

\[
861
\]

Por eso necesitamos expresar incertidumbre.

\end{frame}

% ------------------------------------------------------------

\begin{frame}{Intervalo de confianza}

Un intervalo de confianza proporciona un rango de valores plausibles
para un parámetro bajo un procedimiento estadístico determinado.

\vspace{0.5cm}

De forma aproximada, para ciertos casos:

\[
\bar{x}
\pm
1.96(SE)
\]

para un nivel del 95\%.

\end{frame}

% ------------------------------------------------------------

\begin{frame}{Ejemplo}

Supongamos:

\[
\bar{x}
=
850
\]

y:

\[
SE
=
20
\]

Entonces:

\[
850
\pm
1.96(20)
\]

\[
850
\pm
39.2
\]

Por lo tanto:

\[
[810.8,\ 889.2]
\]

\end{frame}

% ------------------------------------------------------------

\begin{frame}{Interpretación frecuentista}

Una interpretación adecuada es:

\begin{quote}

Si repitiéramos el procedimiento muchas veces,
aproximadamente el 95\% de los intervalos construidos de esta forma
contendrían el parámetro verdadero,
bajo los supuestos del método.

\end{quote}

\vspace{0.4cm}

\begin{block}{Idea práctica}
El intervalo comunica precisión e incertidumbre.
\end{block}

\end{frame}

% ------------------------------------------------------------

\begin{frame}[fragile]{Intervalo de confianza con SciPy}

\begin{lstlisting}
import scipy.stats as stats

datos = df["ventas"].dropna()

media = datos.mean()

sem = stats.sem(datos)

ic = stats.t.interval(
    confidence=0.95,
    df=len(datos)-1,
    loc=media,
    scale=sem
)

print("Media:", media)

print("IC 95%:", ic)
\end{lstlisting}

\end{frame}

% ------------------------------------------------------------

\begin{frame}{Intervalo ancho vs estrecho}

Un intervalo más estrecho suele indicar:

\begin{center}

\Large
mayor precisión

\end{center}

Puede depender de:

\begin{itemize}

    \item tamaño de muestra;

    \item variabilidad;

    \item nivel de confianza;

    \item modelo estadístico.

\end{itemize}

\end{frame}

% ------------------------------------------------------------

\begin{frame}{Precisión no significa ausencia de sesgo}

Podemos tener:

\begin{center}

\Large
un intervalo muy estrecho
\end{center}

pero construido sobre una muestra sesgada.

\vspace{0.5cm}

Entonces tendremos una estimación:

\begin{center}

\Large
\textbf{precisa pero equivocada}
\end{center}

\vspace{0.4cm}

\begin{alertblock}{Lección}
Precisión y exactitud no son exactamente lo mismo.
\end{alertblock}

\end{frame}

% ============================================================
\section{Hipótesis estadísticas}
% ============================================================

\begin{frame}{Una nueva pregunta}

Una empresa modifica su estrategia comercial.

Antes:

\[
\bar{x}_A
=
820
\]

Después:

\[
\bar{x}_B
=
860
\]

\vspace{0.4cm}

\begin{alertblock}{Pregunta}
¿La diferencia refleja un cambio real o podría aparecer por variabilidad muestral?
\end{alertblock}

\end{frame}

% ------------------------------------------------------------

\begin{frame}{Hipótesis nula}

La hipótesis nula se representa como:

\[
H_0
\]

Usualmente representa:

\begin{center}

\Large
ausencia de efecto o diferencia

\end{center}

Por ejemplo:

\[
H_0:
\mu_A
=
\mu_B
\]

\end{frame}

% ------------------------------------------------------------

\begin{frame}{Hipótesis alternativa}

La hipótesis alternativa se representa:

\[
H_1
\]

Ejemplo bilateral:

\[
H_1:
\mu_A
\neq
\mu_B
\]

Ejemplo direccional:

\[
H_1:
\mu_B
>
\mu_A
\]

\end{frame}

% ------------------------------------------------------------

\begin{frame}{La lógica de una prueba estadística}

\begin{center}

\begin{tikzpicture}[
node distance=0.55cm,
every node/.style={
draw,
rounded corners,
fill=azulclaro,
minimum width=4.2cm,
minimum height=0.7cm,
font=\small,
align=center
}
]

\node (h0) {
Suponemos que $H_0$ es cierta
};

\node[below=of h0] (datos) {
Observamos nuestros datos
};

\node[below=of datos] (extremo) {
Medimos qué tan extremo es el resultado
};

\node[below=of extremo] (decision) {
Evaluamos la evidencia
};

\draw[-{Latex}] (h0) -- (datos);

\draw[-{Latex}] (datos) -- (extremo);

\draw[-{Latex}] (extremo) -- (decision);

\end{tikzpicture}

\end{center}

\end{frame}

% ============================================================
\section{Valor p}
% ============================================================

\begin{frame}{¿Qué es el valor p?}

El valor $p$ es:

\begin{quote}

La probabilidad, suponiendo que la hipótesis nula y el modelo del test
son correctos, de obtener un resultado al menos tan extremo como el observado.

\end{quote}

\vspace{0.4cm}

\begin{alertblock}{Muy importante}
El valor $p$ NO es la probabilidad de que $H_0$ sea verdadera.
\end{alertblock}

\end{frame}

% ------------------------------------------------------------

\begin{frame}{Nivel de significancia}

Un valor utilizado frecuentemente es:

\[
\alpha
=
0.05
\]

Si:

\[
p < \alpha
\]

rechazamos $H_0$ bajo ese criterio.

Si:

\[
p \geq \alpha
\]

no encontramos evidencia suficiente para rechazar $H_0$.

\end{frame}

% ------------------------------------------------------------

\begin{frame}{No rechazar no significa demostrar}

Si obtenemos:

\[
p = 0.30
\]

no debemos decir:

\begin{center}

``Demostramos que ambos grupos son iguales.''
\end{center}

\vspace{0.4cm}

La interpretación correcta es:

\begin{center}

``No encontramos evidencia suficiente para rechazar la hipótesis nula.''
\end{center}

\end{frame}

% ============================================================
\section{Errores estadísticos}
% ============================================================

\begin{frame}{Error tipo I}

Ocurre cuando:

\begin{center}

rechazamos $H_0$

\end{center}

cuando en realidad:

\begin{center}

$H_0$ era cierta.

\end{center}

\vspace{0.4cm}

Puede interpretarse como:

\begin{center}

\Large
\textbf{falsa alarma}
\end{center}

\end{frame}

% ------------------------------------------------------------

\begin{frame}{Error tipo II}

Ocurre cuando:

\begin{center}

no rechazamos $H_0$

\end{center}

cuando realmente:

\begin{center}

existe un efecto.

\end{center}

\vspace{0.4cm}

Puede interpretarse como:

\begin{center}

\Large
\textbf{no detectar una señal real}
\end{center}

\end{frame}

% ------------------------------------------------------------

\begin{frame}{Ejemplo empresarial}

\begin{columns}

\column{0.45\textwidth}

\begin{alertblock}{Tipo I}

Concluimos que una campaña funciona.

Pero realmente no produce mejora.

\end{alertblock}

\column{0.45\textwidth}

\begin{exampleblock}{Tipo II}

Concluimos que no existe mejora.

Pero realmente la campaña sí funciona.

\end{exampleblock}

\end{columns}

\end{frame}

% ============================================================
\section{Significancia práctica}
% ============================================================

\begin{frame}{Significancia estadística}

Supongamos:

\begin{center}

\begin{tabular}{lc}
\toprule

Grupo &
Promedio \\

\midrule

A &
\$100.00 \\

B &
\$100.03 \\

\bottomrule

\end{tabular}

\end{center}

Con millones de observaciones podríamos encontrar:

\[
p < 0.001
\]

\begin{alertblock}{Pregunta}
¿Una diferencia de tres centavos importa?
\end{alertblock}

No necesariamente.

\end{frame}

% ------------------------------------------------------------

\begin{frame}{Dos preguntas diferentes}

\begin{columns}

\column{0.45\textwidth}

\begin{block}{Valor p}

Pregunta:

\begin{center}

¿Existe evidencia de una diferencia?

\end{center}

\end{block}

\column{0.45\textwidth}

\begin{exampleblock}{Tamaño del efecto}

Pregunta:

\begin{center}

¿Qué tan grande es la diferencia?

\end{center}

\end{exampleblock}

\end{columns}

\vspace{0.4cm}

\begin{alertblock}{Lección}
Significancia estadística no implica automáticamente relevancia práctica.
\end{alertblock}

\end{frame}

% ============================================================
\section{Live Coding}
% ============================================================

\begin{frame}{Caso práctico: Experimento A/B}

Tenemos dos grupos:

\begin{itemize}

    \item Grupo A: experiencia actual;

    \item Grupo B: nueva experiencia.

\end{itemize}

Queremos determinar:

\begin{center}

\Large
¿El gasto promedio cambió?

\end{center}

Dataset:

\begin{center}

\texttt{experimento\_marketing.csv}

\end{center}

\end{frame}

% ------------------------------------------------------------

\begin{frame}[fragile]{Paso 1: cargar el dataset}

\begin{lstlisting}
import pandas as pd

df = pd.read_csv(
    "experimento_marketing.csv"
)

print(df.head())

print(df.shape)

df.info()
\end{lstlisting}

\end{frame}

% ------------------------------------------------------------

\begin{frame}[fragile]{Paso 2: descriptiva por grupo}

\begin{lstlisting}
resumen = (
    df.groupby("grupo")["gasto"]
      .agg([
          "count",
          "mean",
          "median",
          "std"
      ])
)

print(resumen)
\end{lstlisting}

\begin{block}{Regla}
Antes de aplicar una prueba estadística debemos conocer los datos.
\end{block}

\end{frame}

% ------------------------------------------------------------

\begin{frame}[fragile]{Paso 3: visualizar}

\begin{lstlisting}
import seaborn as sns
import matplotlib.pyplot as plt

sns.boxplot(
    data=df,
    x="grupo",
    y="gasto"
)

plt.title(
    "Distribucion del gasto por grupo"
)

plt.show()
\end{lstlisting}

\end{frame}

% ------------------------------------------------------------

\begin{frame}{Pregunta antes de probar}

Después de observar la gráfica:

\begin{center}

\Large
¿Parece existir una diferencia?
\end{center}

\vspace{0.5cm}

Pero recuerda:

\begin{alertblock}{Regla}
Una diferencia visual no demuestra por sí sola una diferencia estadística.
\end{alertblock}

\end{frame}

% ------------------------------------------------------------

\begin{frame}[fragile]{Paso 4: separar grupos}

\begin{lstlisting}
grupo_a = df.loc[
    df["grupo"] == "A",
    "gasto"
].dropna()

grupo_b = df.loc[
    df["grupo"] == "B",
    "gasto"
].dropna()

print(
    grupo_a.mean()
)

print(
    grupo_b.mean()
)
\end{lstlisting}

\end{frame}

% ------------------------------------------------------------

\begin{frame}{Formular hipótesis}

Hipótesis nula:

\[
H_0:
\mu_A
=
\mu_B
\]

Hipótesis alternativa:

\[
H_1:
\mu_A
\neq
\mu_B
\]

\vspace{0.5cm}

\begin{block}{Pregunta}
¿Existe evidencia de que el gasto promedio de ambos grupos sea diferente?
\end{block}

\end{frame}

% ------------------------------------------------------------

\begin{frame}[fragile]{Paso 5: Welch t-test}

\begin{lstlisting}
from scipy.stats import ttest_ind

t_stat, p_value = ttest_ind(
    grupo_a,
    grupo_b,
    equal_var=False
)

print(
    "t-statistic:",
    t_stat
)

print(
    "p-value:",
    p_value
)
\end{lstlisting}

\end{frame}

% ------------------------------------------------------------

\begin{frame}{¿Por qué Welch?}

La prueba t de Welch permite comparar las medias de dos grupos independientes.

Una ventaja es que:

\begin{center}

no exige asumir exactamente la misma varianza en ambos grupos.

\end{center}

\vspace{0.4cm}

\begin{alertblock}{Importante}
Toda prueba estadística tiene supuestos que debemos considerar.
\end{alertblock}

\end{frame}

% ------------------------------------------------------------

\begin{frame}[fragile]{Paso 6: decisión estadística}

\begin{lstlisting}
alpha = 0.05

if p_value < alpha:

    print(
        "Hay evidencia de una diferencia."
    )

else:

    print(
        "No hay evidencia suficiente."
    )
\end{lstlisting}

\end{frame}

% ------------------------------------------------------------

\begin{frame}{Pero todavía falta algo}

Supongamos:

\[
p = 0.01
\]

Entonces encontramos evidencia estadística.

Pero todavía debemos preguntar:

\begin{center}

\Huge
¿Cuánto cambió?
\end{center}

\vspace{0.4cm}

Ese es el puente hacia la magnitud del efecto.

\end{frame}

% ------------------------------------------------------------

\begin{frame}[fragile]{Paso 7: diferencia observada}

\begin{lstlisting}
media_a = grupo_a.mean()

media_b = grupo_b.mean()

diferencia = (
    media_b
    - media_a
)

print(
    "Media A:",
    media_a
)

print(
    "Media B:",
    media_b
)

print(
    "Diferencia:",
    diferencia
)
\end{lstlisting}

\end{frame}

% ============================================================
\section{Tamaño del efecto}
% ============================================================

\begin{frame}{Tamaño del efecto}

Una medida clásica es Cohen's $d$:

\[
d
=
\frac{
\bar{x}_1
-
\bar{x}_2
}{
s_{\text{pooled}}
}
\]

\vspace{0.5cm}

Conceptualmente:

\begin{center}

\textbf{Valor p}

$\downarrow$

¿Existe evidencia?

\vspace{0.5cm}

\textbf{Effect Size}

$\downarrow$

¿Qué tan grande es el efecto?

\end{center}

\end{frame}

% ============================================================
\section{Correlación}
% ============================================================

\begin{frame}[fragile]{Correlación}

Podemos estudiar asociaciones entre variables numéricas:

\begin{lstlisting}
correlacion = df[
    [
        "publicidad",
        "ventas"
    ]
].corr()

print(correlacion)
\end{lstlisting}

Podríamos obtener:

\[
r
=
0.82
\]

\end{frame}

% ------------------------------------------------------------

\begin{frame}{Interpretación de correlación}

Un valor cercano a:

\[
+1
\]

indica una asociación lineal positiva fuerte.

Un valor cercano a:

\[
-1
\]

indica una asociación lineal negativa fuerte.

Un valor cercano a:

\[
0
\]

indica poca asociación lineal.

\end{frame}

% ------------------------------------------------------------

\begin{frame}{Correlación no implica causalidad}

Encontrar:

\[
X
\leftrightarrow
Y
\]

no demuestra:

\[
X
\rightarrow
Y
\]

Podrían existir:

\begin{itemize}

    \item variables confusoras;

    \item causalidad inversa;

    \item selección;

    \item coincidencias;

    \item relaciones indirectas.

\end{itemize}

\end{frame}

% ------------------------------------------------------------

\begin{frame}{Ejemplo clásico}

Podemos observar:

\begin{center}

Ventas de helado

$\uparrow$

\vspace{0.3cm}

Casos de insolación

$\uparrow$

\end{center}

¿El helado causa insolación?

\vspace{0.5cm}

No necesariamente.

Una tercera variable podría explicar ambas:

\begin{center}

\Large
Temperatura
\end{center}

\end{frame}

% ============================================================
\section{Statistical Detective Challenge}
% ============================================================

\begin{frame}{Misión}

\begin{center}

{\Huge
\textcolor{azulmarista}{
\textbf{Statistical Detective}
}
}

\vspace{0.7cm}

Marketing afirma:

\vspace{0.3cm}

{\Large
\textbf{``La campaña B es mejor que la campaña A.''}
}

\vspace{0.8cm}

Su misión:

\vspace{0.2cm}

\textbf{determinar si los datos respaldan esa afirmación.}

\end{center}

\end{frame}

% ------------------------------------------------------------

\begin{frame}{Dataset del reto}

Archivo:

\begin{center}

\texttt{campania\_ab.csv}

\end{center}

Variables:

\begin{itemize}

    \item \texttt{id\_usuario};

    \item \texttt{grupo};

    \item \texttt{edad};

    \item \texttt{region};

    \item \texttt{conversion};

    \item \texttt{gasto};

    \item \texttt{tiempo\_sitio}.

\end{itemize}

\end{frame}

% ------------------------------------------------------------

\begin{frame}{Paso 1: explorar}

Cada equipo debe determinar:

\begin{itemize}

    \item tamaño de los grupos;

    \item media;

    \item mediana;

    \item desviación estándar;

    \item valores faltantes;

    \item posibles outliers.

\end{itemize}

\end{frame}

% ------------------------------------------------------------

\begin{frame}[fragile]{Código de apoyo}

\begin{lstlisting}
resumen = (
    df.groupby("grupo")["gasto"]
      .agg([
          "count",
          "mean",
          "median",
          "std"
      ])
)

print(resumen)
\end{lstlisting}

\end{frame}

% ------------------------------------------------------------

\begin{frame}{Paso 2: visualizar}

Cada equipo debe crear al menos una gráfica.

Opciones:

\begin{itemize}

    \item boxplot;

    \item histograma;

    \item gráfica de comparación de medias.

\end{itemize}

\vspace{0.4cm}

\begin{block}{Regla}
La visualización debe ayudar a responder la pregunta.
\end{block}

\end{frame}

% ------------------------------------------------------------

\begin{frame}{Paso 3: formular hipótesis}

Para comparar gasto promedio:

\[
H_0:
\mu_A
=
\mu_B
\]

\[
H_1:
\mu_A
\neq
\mu_B
\]

\vspace{0.5cm}

Deben explicar estas hipótesis usando lenguaje de negocio.

\end{frame}

% ------------------------------------------------------------

\begin{frame}{Paso 4: realizar la prueba}

Para este ejercicio:

\begin{center}

\Large
Welch t-test

\end{center}

Deben reportar:

\begin{itemize}

    \item estadístico t;

    \item valor $p$;

    \item nivel $\alpha$;

    \item decisión estadística.

\end{itemize}

\end{frame}

% ------------------------------------------------------------

\begin{frame}{Paso 5: medir la diferencia}

No basta con decir:

\begin{center}

\texttt{p = 0.03}

\end{center}

También deben calcular:

\[
\bar{x}_B
-
\bar{x}_A
\]

y responder:

\begin{center}

\Large
¿La diferencia es relevante?
\end{center}

\end{frame}

% ------------------------------------------------------------

\begin{frame}{Trampas del dataset}

El dataset puede incluir:

\begin{itemize}

    \item tamaños de muestra diferentes;

    \item valores extremos;

    \item distribuciones sesgadas;

    \item varianzas diferentes;

    \item pequeñas diferencias estadísticamente significativas.

\end{itemize}

\vspace{0.4cm}

\begin{alertblock}{Objetivo}
No ejecutar automáticamente una prueba.
Primero entender los datos.
\end{alertblock}

\end{frame}

% ------------------------------------------------------------

\begin{frame}{Formato de entrega}

Cada equipo deberá presentar:

\begin{enumerate}

    \item pregunta;

    \item hipótesis;

    \item resumen descriptivo;

    \item visualización;

    \item prueba estadística;

    \item valor $p$;

    \item diferencia observada;

    \item interpretación;

    \item recomendación.

\end{enumerate}

\end{frame}

% ============================================================
\section{Errores comunes}
% ============================================================

\begin{frame}{Error común 1}

\begin{center}

\Large
``p = 0.03, entonces existe 97\% de probabilidad de que mi hipótesis sea cierta.''

\end{center}

\vspace{0.6cm}

\begin{alertblock}{Incorrecto}
El valor $p$ no es la probabilidad de que una hipótesis sea verdadera.
\end{alertblock}

\end{frame}

% ------------------------------------------------------------

\begin{frame}{Error común 2}

\begin{center}

\Large
``p > 0.05, entonces demostramos que los grupos son iguales.''

\end{center}

\vspace{0.6cm}

\begin{alertblock}{Incorrecto}
No encontrar evidencia suficiente de diferencia no demuestra igualdad.
\end{alertblock}

\end{frame}

% ------------------------------------------------------------

\begin{frame}{Error común 3}

\begin{center}

\Large
``p < 0.05, entonces el resultado es importante.''

\end{center}

\vspace{0.6cm}

\begin{alertblock}{Incorrecto}
Significancia estadística no implica necesariamente importancia práctica.
\end{alertblock}

\end{frame}

% ------------------------------------------------------------

\begin{frame}{Error común 4}

\begin{center}

\Large
``Tengo muchos datos, entonces mi análisis debe ser correcto.''

\end{center}

\vspace{0.6cm}

\begin{alertblock}{Incorrecto}
Un dataset enorme puede seguir teniendo sesgos o problemas de diseño.
\end{alertblock}

\end{frame}

% ============================================================
\section{Pipeline estadístico}
% ============================================================

\begin{frame}{Pipeline de análisis estadístico}

\begin{center}

\begin{tikzpicture}[
node distance=0.40cm,
every node/.style={
draw,
rounded corners,
fill=azulclaro,
minimum width=3.3cm,
minimum height=0.50cm,
font=\scriptsize
}
]

\node (pregunta) {Pregunta};

\node[below=of pregunta] (explorar) {Explorar};

\node[below=of explorar] (describir) {Describir};

\node[below=of describir] (hipotesis) {Formular hipótesis};

\node[below=of hipotesis] (test) {Aplicar prueba};

\node[below=of test] (magnitud) {Medir magnitud};

\node[below=of magnitud] (interpretar) {Interpretar};

\draw[-{Latex}] (pregunta) -- (explorar);

\draw[-{Latex}] (explorar) -- (describir);

\draw[-{Latex}] (describir) -- (hipotesis);

\draw[-{Latex}] (hipotesis) -- (test);

\draw[-{Latex}] (test) -- (magnitud);

\draw[-{Latex}] (magnitud) -- (interpretar);

\end{tikzpicture}

\end{center}

\end{frame}

% ============================================================
\section{Herramientas}
% ============================================================

\begin{frame}{Herramientas de Pandas}

\small

\begin{table}
\centering

\begin{tabular}{ll}
\toprule

\textbf{Objetivo} &
\textbf{Función} \\

\midrule

Media &
\texttt{mean()} \\

Mediana &
\texttt{median()} \\

Moda &
\texttt{mode()} \\

Varianza &
\texttt{var()} \\

Desviación &
\texttt{std()} \\

Percentiles &
\texttt{quantile()} \\

Muestreo &
\texttt{sample()} \\

Agrupar &
\texttt{groupby()} \\

Resumen múltiple &
\texttt{agg()} \\

\bottomrule

\end{tabular}

\end{table}

\end{frame}

% ------------------------------------------------------------

\begin{frame}{Herramientas de SciPy}

\small

\begin{table}
\centering

\begin{tabular}{ll}
\toprule

\textbf{Objetivo} &
\textbf{Herramienta} \\

\midrule

Error estándar &
\texttt{stats.sem()} \\

Intervalo de confianza &
\texttt{stats.t.interval()} \\

Comparación de medias &
\texttt{ttest\_ind()} \\

Z-score &
\texttt{zscore()} \\

\bottomrule

\end{tabular}

\end{table}

\end{frame}

% ============================================================
\section{Ideas clave}
% ============================================================

\begin{frame}{12 ideas que deben recordar}

\scriptsize

\begin{enumerate}

    \item Una muestra es una ventana hacia una población.

    \item Una muestra grande no elimina automáticamente el sesgo.

    \item La media no siempre representa correctamente los datos.

    \item La dispersión importa tanto como el centro.

    \item Las distribuciones contienen información importante.

    \item Los estadísticos cambian entre muestras.

    \item El Teorema Central del Límite explica el comportamiento de medias muestrales.

    \item El error estándar mide incertidumbre de una estimación.

    \item Los intervalos de confianza comunican precisión.

    \item El valor $p$ no es la probabilidad de que $H_0$ sea verdadera.

    \item Significancia estadística no implica importancia práctica.

    \item Correlación no implica causalidad.

\end{enumerate}

\end{frame}

% ============================================================
\section{Volvamos al problema}
% ============================================================

\begin{frame}{Volvamos al inicio}

Al comenzar preguntamos:

\begin{center}

A:

\[
12.1\%
\]

B:

\[
13.0\%
\]

\vspace{0.4cm}

{\Huge
¿B es mejor?
}

\end{center}

\end{frame}

% ------------------------------------------------------------

\begin{frame}{Ahora nuestra respuesta cambia}

No podemos responder solamente comparando:

\[
12.1\%
\]

contra:

\[
13.0\%
\]

Necesitamos considerar:

\begin{itemize}

    \item tamaño de muestra;

    \item variabilidad;

    \item incertidumbre;

    \item diseño del experimento;

    \item significancia estadística;

    \item tamaño del efecto;

    \item relevancia para el negocio.

\end{itemize}

\end{frame}

% ============================================================
\section{Puente a Series Temporales}
% ============================================================

\begin{frame}{¿Qué cambia cuando aparece el tiempo?}

Hasta ahora hemos analizado observaciones sin centrarnos demasiado
en el momento en que ocurrieron.

Pero consideremos:

\begin{center}

Ventas de enero

$\downarrow$

Ventas de febrero

$\downarrow$

Ventas de marzo

$\downarrow$

Ventas de abril

\end{center}

\vspace{0.4cm}

Aquí:

\begin{center}

\Large
\textbf{el orden importa}
\end{center}

\end{frame}

% ------------------------------------------------------------

\begin{frame}{Una nueva clase de datos}

Supongamos que tenemos:

\[
X_1,
X_2,
X_3,
\ldots,
X_t
\]

donde cada observación corresponde a un momento diferente.

\vspace{0.4cm}

Aparecen preguntas nuevas:

\begin{itemize}

    \item ¿Los valores actuales dependen del pasado?

    \item ¿Existe una tendencia?

    \item ¿Hay patrones que se repiten?

    \item ¿Existen ciclos?

    \item ¿Existe estacionalidad?

    \item ¿Podemos utilizar el pasado para entender el futuro?

\end{itemize}

\end{frame}

% ------------------------------------------------------------

\begin{frame}{Próxima clase}

\begin{center}

{\Huge
\textbf{Clase 7}
}

\vspace{0.4cm}

{\LARGE
\textcolor{azulmarista}{
\textbf{Series Temporales I}
}
}

\vspace{0.3cm}

{\Large
Tendencia, estacionalidad y estructura temporal
}

\end{center}

\vspace{0.5cm}

Analizaremos:

\begin{itemize}

    \item índice temporal;

    \item resampling;

    \item tendencia;

    \item estacionalidad;

    \item ciclos;

    \item ruido;

    \item medias móviles;

    \item ventanas móviles;

    \item diferenciación;

    \item autocorrelación;

    \item ACF y PACF;

    \item estacionariedad.

\end{itemize}

\end{frame}

% ============================================================
% CIERRE
% ============================================================

\begin{frame}

\begin{center}

{\Huge
\textcolor{azulmarista}{
\textbf{Estadística para Data Science}
}
}

\vspace{0.8cm}

{\Large
Los datos no hablan por sí solos.
}

\vspace{0.4cm}

{\Large
La estadística nos ayuda a distinguir
}

\vspace{0.3cm}

{\LARGE
\textbf{señal}
}

\vspace{0.1cm}

de

\vspace{0.1cm}

{\LARGE
\textbf{ruido}.
}

\vspace{0.8cm}

\textbf{Diplomado de Python y Análisis de Datos}

\vspace{0.2cm}

\textbf{Universidad Marista}

\end{center}

\end{frame}

% ============================================================

\end{document}